\documentclass[10pt,a4paper]{amsart}
\usepackage[margin=19mm]{geometry}
\usepackage{amsmath,amssymb,mathtools}
\usepackage[scr=rsfs,cal=euler]{mathalfa}
\usepackage{xfrac}
\usepackage[unicode,hidelinks,bookmarks=false]{hyperref}
\newcommand{\ie}{i.e.\ }
\newcommand{\cf}{cf.\ }
\newcommand{\resp}{respectively\ }
\newcommand{\Subsection}{\subsection}
\providecommand{\nobreakdash}{\nobreak\hbox{-}}
\setlength{\emergencystretch}{2em}
\multlinegap0pt
\usepackage{amssymb}
\newtheorem{lemma}{Lemma}
\newcommand{\RR}{\mathbb R}
\newcommand{\ZZ}{\mathbb Z}
\newcommand{\dint}{\,\mathrm d}
\newcommand{\e}{\mathrm e}
\newcommand{\eps}{\varepsilon}
\newcommand{\h}{\mathbf h}
\newcommand{\x}{\mathbf{x}}
\title{Improved bounds in Birch's theorem for forms in many variables}
\author{Yijie Diao, Lena Wurzinger}
\date{Source: arXiv:2609.00757v1 (CC BY 4.0)}
\begin{document}
\maketitle

\noindent\emph{Source and licence.} Improved bounds in Birch's theorem for forms in many variables. Version arXiv:2609.00757v1 (CC BY 4.0), distributed by arXiv under the Creative Commons Attribution 4.0 International licence, http://creativecommons.org/licenses/by/4.0/, as recorded in the arXiv licence metadata for this exact version and shown on the abstract page https://arxiv.org/abs/2609.00757v1. Full licence notice: \texttt{licenses/arxiv-2609.00757-CC-BY-4.0.txt}.\par\smallskip

\noindent\textit{Editorial scope.} Counted unit: the article's lemma lem:average (source0.tex lines 354--370). The article's complete original proof of it is delivered with it (source0.tex lines 372--502, one \texttt{proof} environment of 131 lines). No mathematics is written, repaired or re-derived: the statement and the proof are the article's own lines, in the article's own notation, and no retained line is substituted or re-flowed.  Retained uncounted as the source-local context the proof needs: lines 331--334; lines 337--340; lines 342--344; lines 346--350.  Nothing was omitted from any retained span, so the package carries no omission record.  No retained line contains a citation, so the wrapper carries no bibliography and no bibliography entry.  No retained line contains a figure environment, an image-inclusion command or drawing code, so no image is delivered and the package declares no asset dependency.  Editorial formatting only, in addition: an AMS \texttt{amsart} preamble that restates the article's own declarations and macros used by the retained lines, namely the environments \texttt{lemma} on separate counters, together with the standard packages those lines need.  The retained environments print in this document as Lemma 1 in order of appearance, whereas the article prints the same items with the numbering of its own full document.  The complete native source of the article as distributed in the arXiv e-print for this exact version is bundled unchanged as \texttt{sources/209059/source0.tex}.
\begin{equation}\label{eq:lower-bound-partial-F}
 |\partial_1F(\x)|\geq c_1B^{d-1},
 \qquad\text{for }\x/B\in\operatorname{supp}(\omega).
\end{equation}

\begin{equation}\label{eq:exponential-sum}
 S_\omega(\alpha;B)=\sum_{\x\in\ZZ^n}\omega(\x/B)
 \e(\alpha F(\x)).
\end{equation}

\begin{equation}\label{eq:first-difference}
 \Delta_\h F(\x)=F(\x+\h)-F(\x).
\end{equation}

\begin{equation}\label{eq:first-difference-sum}
 S_{\h}(\alpha;B)=\sum_{\x\in\ZZ^n}
 \omega((\x+\h)/B)\omega(\x/B)
 \e(\alpha\Delta_\h F(\x)).
\end{equation}

\begin{lemma}\label{lem:average}
Let \(\eps>0\), \(B\geq2\), \(1\leq H\leq B\), and \(0<t\leq1\).
Let \(a\in\ZZ\) and \(q\in\mathbb N\) with \((a,q)=1\), and put \(\lambda=(HB^{d-1})^{-1}\).
Define
\begin{equation}\label{eq:shift-average}
 \Sigma_H=\max_{(t-B^\eps\lambda)_+\leq |z|
 \leq2t+B^\eps\lambda}
 \sum_{|\h|\leq HB^\eps}|S_{\h}(a/q+z;B)|.
\end{equation}
Then, for every fixed \(J\geq1\), we have
\begin{equation}\label{eq:average}
 \int_{t\leq|z|\leq2t}|S_\omega(a/q+z;B)|\dint z
 \ll B^{-J}+B^\eps(t+\lambda)
 \left(\frac BH\right)^{(n-1)/2}\Sigma_H^{1/2}.
\end{equation}
The implied constant depends only on \(F,\omega,\eps\), and \(J\).
\end{lemma}

\begin{proof}
Let \(J\) be fixed.
Cover \(t\leq|z|\leq2t\) by \(O(1+t/\lambda)\) intervals \(I(z_0)\) of radius \(\lambda\), choosing every center \(z_0\) in the set \(\{z:t\leq|z|\leq2t\}\).
For each center \(z_0\), use \eqref{eq:exponential-sum} to put
\[
 \mathcal M(z_0)=\int_{\RR}
 \exp(-(z-z_0)^2/\lambda^2)
 |S_\omega(a/q+z;B)|^2\dint z.
\]
The Gaussian is \(\gg 1\) on \(I(z_0)\), so Cauchy--Schwarz gives
\begin{equation}\label{eq:interval-cauchy}
 \int_{I(z_0)}|S_\omega(a/q+z;B)|\dint z
 \ll\lambda^{1/2}\mathcal M(z_0)^{1/2}.
\end{equation}

We first record the van der Corput inequality for lattice sums in the form needed here.
If \(\mathcal H\subset\ZZ^n\) is finite and
\[
 \#\{\mathbf y\in\ZZ^n:f(\mathbf y+\mathbf u)\neq0
 \text{ for some }\mathbf u\in\mathcal H\}\ll B^n,
\]
then translation averaging followed by Cauchy--Schwarz gives
\begin{equation}\label{eq:general-vdc}
 \begin{aligned}
 (\#\mathcal H)^2\left|\sum_{\x\in\ZZ^n}f(\x)\right|^2
 &\ll B^n\sum_{\mathbf y\in\ZZ^n}
 \left|\sum_{\mathbf u\in\mathcal H}f(\mathbf y+\mathbf u)\right|^2\\
 &=B^n\sum_{\h\in\ZZ^n}N(\h)
 \sum_{\mathbf y\in\ZZ^n}
 f(\mathbf y+\h)\overline{f(\mathbf y)}.
 \end{aligned}
\end{equation}
where
\[
 N(\h)=\#\{(\mathbf u,\mathbf v)\in\mathcal H^2: \mathbf u-\mathbf v=\h\}.
\]

Choose \(\delta>0\) sufficiently small, depending only on \(F\) and \(\omega\).
Apply \eqref{eq:general-vdc} to
\[
 f_z(\x)=\omega(\x/B)\e((a/q+z)F(\x)),
\]
with
\[
\mathcal H=
 \bigl([1,\delta B]\times[1,H]^{n-1}\bigr)\cap\ZZ^n.
\]
Since \(\omega\) has fixed compact support and \(|\mathbf u|\ll B\) for \(\mathbf u\in\mathcal H\), this support condition holds.
After enlarging the implied constant we may assume that \(B\) is large enough that \(\#\mathcal H\asymp BH^{n-1}\); the bounded remaining values of \(B\) are harmless.
Moreover, \(N(\h)\leq\#\mathcal H\), and the correlation in \eqref{eq:general-vdc} is exactly \(S_\h(a/q+z;B)\).
The difference set is contained in \([ -\delta B,\delta B]\times[-H,H]^{n-1}\).
Multiply \eqref{eq:general-vdc} by the Gaussian and integrate in \(z\).
Using the triangle inequality and the bound on \(N(\h)\), we obtain
\begin{equation}\label{eq:gaussian-vdc}
 \mathcal M(z_0)\ll
 \left(\frac BH\right)^{n-1}\!
 \sum_{\substack{|h_1|\leq\delta B\\
                  |h_j|\leq H\\
                  2\leq j\leq n}}
 \left|\int_{\RR}\exp(-(z-z_0)^2/\lambda^2)
 S_{\h}(a/q+z;B)\dint z\right|.
\end{equation}
For \(u\in\RR\), we have the identity
\begin{equation*}
 \int_{\RR}\exp(-(z-z_0)^2/\lambda^2)
 \e((a/q+z)u)\dint z
 =\lambda\sqrt\pi
 \e((a/q+z_0)u)
 \exp(-\pi^2\lambda^2u^2).
\end{equation*}
Apply this with \(u=\Delta_\h F(\x)\) after expanding \(S_{\h}\) in \eqref{eq:first-difference-sum}.
The summand vanishes unless \(\omega(\x/B)\omega((\x+\h)/B)\neq0\), and on this support Taylor's formula applied to \eqref{eq:first-difference} gives
\begin{equation}\label{eq:first-difference-taylor}
 \Delta_\h F(\x)
 =h_1\partial_1F(\x)+O(HB^{d-1} + (h_1^2+H^2)B^{d-2}).
\end{equation}
For every shift in \eqref{eq:gaussian-vdc} with \(|h_1|>HB^\eps\), we have
\begin{equation*}
 HB^{d-1}+(h_1^2+H^2)B^{d-2}
 \ll(\delta+B^{-\eps})|h_1|B^{d-1}.
\end{equation*}
Choose \(\delta\) sufficiently small in terms of the constant in \eqref{eq:first-difference-taylor} and the lower bound \eqref{eq:lower-bound-partial-F}.
For all sufficiently large \(B\), the error term is then at most half the main term, and hence
\[
 |\Delta_\h F(\x)|\geq \frac{c_1}{2}|h_1|B^{d-1},
\]
for every such shift.
The remaining bounded values of \(B\) are absorbed into the final implied constant.
Hence the factor in the preceding Fourier transform satisfies
\[
 \exp(-\pi^2\lambda^2(\Delta_\h F(\x))^2)
 \leq\exp(-c\lambda^2|h_1|^2B^{2d-2})
 =\exp(-c|h_1|^2/H^2)
 \leq\exp(-cB^{2\eps}).
\]
Here \(c=\pi^2c_1^2/4>0\).
Choose an auxiliary exponent \(J_0=J_0(J,d,n)\) large enough that the subsequent polynomial losses leave an error \(O(B^{-J})\).
Since the sums over \(\x\) and \(\h\) have polynomial length, this bound shows that the shifts with \(|h_1|>HB^\eps\) contribute \(O_{J_0}(B^{-J_0})\).
The Gaussian weight in \eqref{eq:gaussian-vdc}, together with the trivial estimate \(|S_{\h}(a/q+z;B)|\ll B^n\), similarly shows that the range \(|z-z_0|>B^\eps\lambda\) contributes \(O_{J_0}(B^{-J_0})\).
For the remaining shifts and frequencies, the triangle inequality gives
\begin{equation*}
 \begin{aligned}
 &\sum_{\substack{|h_1|\leq HB^\eps\\
                  |h_j|\leq H\\
                  2\leq j\leq n}}
 \left|\int_{|z-z_0|\leq B^\eps\lambda}
 \exp(-(z-z_0)^2/\lambda^2)
 S_{\h}(a/q+z;B)\dint z\right|\\
 &\qquad\leq
 \int_{|z-z_0|\leq B^\eps\lambda}
 \exp(-(z-z_0)^2/\lambda^2)
 \sum_{|\h|\leq HB^\eps}|S_{\h}(a/q+z;B)|\dint z\\
 &\qquad\ll B^\eps\lambda\Sigma_H.
 \end{aligned}
\end{equation*}
Combining with \eqref{eq:gaussian-vdc} and the two tail estimates gives
\begin{equation}\label{eq:local-average}
 \mathcal M(z_0)
 \ll_{J_0} B^{-J_0}
 +B^\eps\lambda\left(\frac BH\right)^{n-1}\Sigma_H.
\end{equation}
Because every center lies in the original annulus, the truncated frequencies in the preceding display lie in the range used to define \(\Sigma_H\).
Finally, the cover has \(O(1+t/\lambda)\) intervals.
The main term obtained from \eqref{eq:interval-cauchy} and \eqref{eq:local-average} is
\[
 \ll B^\eps(t+\lambda)
 \left(\frac BH\right)^{(n-1)/2}\Sigma_H^{1/2}.
\]
Since \(t\leq1\) and \(\lambda^{-1/2}\leq B^{d/2}\), choosing \(J_0\) sufficiently large makes the summed error \(O(B^{-J})\).
This proves \eqref{eq:average}.
\end{proof}

\end{document}
